% Part: proof-theory
% Chapter: proof-search
% Section: tableaux

\documentclass[../../../include/open-logic-section]{subfiles}

\begin{document}

\olfileid{pt}{ps}{tab}

\olsection{تابلوها}

دستگاه‌های تابلو، دستگاه‌های !!{proof}ی‌اند که ارتباط نزدیکی با
حساب‌های دنباله‌ای دارند و به‌ویژه برای جست‌وجوی برهان، چه دستی و چه
پیاده‌سازی‌شده در نرم‌افزار، مناسب‌اند. در تابلوها به جای ساختن
درخت‌هایی از دنباله‌ها، !!a{proof} درختی از !!{formula}s است. اما
برخلاف استنتاج طبیعی، قواعد آن‌ها بازتاب قواعد حساب دنباله‌ای‌اند. یک
!!{proof} تابلویی در اصل جست‌وجویی صریح برای یک مدل است. ازاین‌رو
گاه آن‌ها را \emph{تابلوهای معنایی} یا \emph{درخت‌های صدق} می‌نامند.

یک تابلوی \emph{نشاندار} درختی از !!{formula}s نشاندار است که
می‌توانند به شکل $\sFmla{\True}{!A}$ یا $\sFmla{\False}{!A}$ باشند.
درخت‌ها رو به پایین رشد می‌کنند. آن‌ها با !!{formula}s آغاز می‌شوند
که آنچه را می‌خواهیم !!{prove} کنیم مشخص می‌سازند. برای نمونه، برای
!!{prove} دنبالهٔ $!A, !B \Sequent !C, !D$، تابلو با
$\sFmla{\True}{!A}, \sFmla{\True}{!B},
\sFmla{\False}{!C}, \sFmla{\False}{!D}$ آغاز می‌شود.
!!^a{structure} که $!A$ و $!B$ را صادق و $!C$ و $!D$ را کاذب سازد،
!!a{structure}ی است که نشان می‌دهد
$!A, !B \Sequent !C, !D$ معتبر نیست.

گره برگ (پایین‌ترین گره) تابلو را می‌توان با قواعد گسترش شاخه توسعه
داد؛ این قواعد یا گره‌های دیگری به همان شاخه می‌افزایند، یا دو گرهٔ
هم‌نیا را در پایین می‌افزایند و شاخه را به دو شاخهٔ تازه تقسیم
می‌کنند. قواعد را می‌توان بر هر !!{formula} نشاندار بالای گره در همان
شاخه اعمال کرد. قواعد گسترش شاخه برای دستگاه تابلویی کلاسیک \Log{Tc}
در \olref{tab:Tc} آمده‌اند.

\subfile{rules-Tc}

همان‌گونه که می‌بینید، این قواعد همان قواعد \Log{G3c} وارونه‌اند و
نشان‌های $\True$ و $\False$ مشخص می‌کنند که !!a{formula} به‌ترتیب در
سمت چپ یا راست یک دنباله فرمول اصلی است.

یک شاخهٔ تابلو \emph{بسته} است اگر شامل $\sFmla{\True}{\lfalse}$
باشد، یا هم $\sFmla{\True}{!A}$ و هم $\sFmla{\False}{!A}$ را در بر
داشته باشد. اگر همهٔ شاخه‌ها بسته باشند، کل تابلو را بسته می‌نامیم.
برای نمونه، تابلوی بستهٔ زیر برای
$!A \lor (!B \land !C) \Sequent (!A \lor !B) \land (!A \lor !C)$
است؛ یعنی با $\sFmla{\True}{!A \lor (!B \land !C)}$ و
$\sFmla{\False}{(!A \lor !B) \land (!A \lor !C)}$ آغاز می‌شود:
\begin{oltableau}
  [\sFmla{\True}{\formula{A} \lor (\formula{B} \land \formula{C})},
    just=\TAss, checked
    [\sFmla{\False}{(\formula{A} \lor \formula{B}) \land
        (\formula{A} \lor \formula{C})}, just=\TAss, checked
      [\sFmla{\True}{\formula{A}}, just={\TRule{\True}{\lor}[1]}
        [\sFmla{\False}{\formula{A} \lor \formula{B}},
          just={\TRule{\False}{\land}[2]}, checked
          [\sFmla{\False}{\formula{A}},
            just={\TRule{\False}{\lor}[4]}
            [\sFmla{\False}{\formula{B}},
              just={\TRule{\False}{\lor}[4]}, close]
          ]
        ]
        [\sFmla{\False}{\formula{A} \lor \formula{C}},
          just={\TRule{\False}{\land}[2]}, checked
          [\sFmla{\False}{\formula{A}},
            just={\TRule{\False}{\lor}[4]}
            [\sFmla{\False}{\formula{C}},
              just={\TRule{\False}{\lor}[4]}, close]
          ]
        ]
      ]
      [\sFmla{\True}{\formula{B} \land \formula{C}},
        just={\TRule{\True}{\lor}[1]}, checked
        [\sFmla{\False}{\formula{A} \lor \formula{B}},
          just={\TRule{\False}{\land}[2]}, checked
          [\sFmla{\True}{\formula{B}},
            just={\TRule{\True}{\land}[3]}, move by=2
            [\sFmla{\True}{\formula{C}},
              just={\TRule{\True}{\land}[3]}
              [\sFmla{\False}{\formula{A}},
                just={\TRule{\False}{\lor}[4]}
                [\sFmla{\False}{\formula{B}},
                  just={\TRule{\False}{\lor}[4]},close
                ]
              ]
            ]
          ]
        ]
        [\sFmla{\False}{\formula{A} \lor \formula{C}},
          just={\TRule{\False}{\land}[2]}, checked
          [\sFmla{\True}{\formula{B}},
            just={\TRule{\True}{\land}[3]}, move by=2
              [\sFmla{\True}{\formula{C}},
                just={\TRule{\True}{\land}[3]}
                [\sFmla{\False}{\formula{A}},
                  just={\TRule{\False}{\lor}[4]}
                  [\sFmla{\False}{\formula{C}},
                  just={\TRule{\False}{\lor}[4]},close
                ]
              ]
            ]
          ]
        ]
      ]
    ]
  ]
\end{oltableau}

علامت‌های تیک نشان می‌دهند قاعدهٔ متناظر بر !!a{formula} در همهٔ
شاخه‌های زیر آن اعمال شده است. نماد $\otimes$ بسته‌بودن شاخه را نشان
می‌دهد. شمارهٔ خط پس از هر قاعده مشخص می‌کند قاعده بر کدام
!!{formula}s نشاندار اعمال شده است (یعنی !!{formula} همان شاخه در خط
مشخص‌شده).

جست‌وجوی برهان در تابلوها آسان است. تابلو را مرحله‌به‌مرحله تولید
می‌کنیم. در مرحلهٔ~$0$ تنها !!{formula}s آغازین را در بالا می‌نویسیم.
در مرحلهٔ~$n+1$ برای هر گره برگ روی یک شاخهٔ باز چنین می‌کنیم: نخستین
(بالاترین) فرمول نشاندار غیراتمی را که هنوز تیک نخورده است بیابید و
قاعدهٔ گسترش شاخه را بر آن اعمال کنید. وقتی قاعده در هر شاخهٔ بازی که
آن فرمول نشاندار را شامل می‌شود اعمال شد، بر آن تیک بزنید. بنابراین
فقط فرمول‌های غیراتمی در شاخه‌های باز که نباید تکرار شوند تیک
می‌خورند. (مثال بالا با همین الگوریتم تولید شده است.)

!!{formula}s نشاندار به شکل $\sFmla{\True}{\lforall}$ و
$\sFmla{\False}{\lexists}$ هرگز تیک نمی‌خورند و باید جداگانه بررسی
شوند. اگر حاضر باشند، باید تضمین کنیم قواعد
$\TRule{\True}{\lforall}$ و $\TRule{\False}{\lexists}$ سرانجام برای
همهٔ ترم‌های قابل اعمال به کار روند. برای نمونه، می‌توان در هر مرحله
آن‌ها را بر همهٔ چنین فرمول‌هایی در هر شاخه اعمال کرد، پس از رسیدگی به
بالاترین !!{formula} نشانداری که از این دو صورت نیست. ترم~$t$ نخستین
ترمی است که می‌توان آن را از !!{constant}s از پیش رخ‌داده در شاخه
(یا $\Obj c_0$ اگر هیچ !!{constant}ی رخ نداده باشد) و
!!{function}s رخ‌داده در !!{formula}s آغازین ساخت، به‌گونه‌ای که
!!{formula} نشاندار متناظر $\sFmla{\True}{!B(t)}$ یا
~$\sFmla{\False}{!B(t)}$ از پیش در شاخه رخ نداده باشد. اگر برای یک
قاعده همهٔ ترم‌های قابل ساخت از پیش نمونه‌ای در شاخه داشته باشند، آن
قاعده در آن مرحله گره تازه‌ای نمی‌افزاید.

اگر این جست‌وجوی برهان تابلوی بسته‌ای تولید نکند، یا شاخهٔ باز متناهی‌ای
دارد که همهٔ !!{formula}s غیراتمی آن بررسی شده‌اند، یا جست‌وجو
درختی نامتناهی و متناهی‌شاخه تولید می‌کند که باید شاخه‌ای نامتناهی
داشته باشد. درست مانند \olref[cpl]{sec}، می‌توانیم
!!a{structure}~$\Struct{M}$ را چنان تعریف کنیم که اگر
$\sFmla{\True}{!A}$ در شاخه رخ دهد، $\Sat{M}{!A}$؛ و اگر
$\sFmla{\False}{!A}$ رخ دهد، $\Sat/{M}{!A}$. (بگذارید
$\Theta = \Setabs{!A}{\sFmla{\True}{!A} \text{ در شاخه}}$ و
$\Xi = \Setabs{!A}{\sFmla{\False}{!A} \text{ در شاخه}}$.)

تابلوهایی که با این روش جست‌وجوی برهان تولید می‌شوند را می‌توان به‌آسانی
با اختصاص دنباله‌ای به هر گره، به برهان‌های \Log{G3c} ترجمه کرد. اگر
فرمول‌های آغازین متناظر با دنبالهٔ $\Gamma \Sequent \Delta$ باشند،
هر یک از فرمول‌های آغازین را با $\Gamma \Sequent \Delta$ برچسب
بزنید. اگر شاخه با اعمال قاعده‌ای گسترشی بر !!{formula} نشاندار
$\sFmla{\True}{!A}$ توسعه یابد، برگ با
$!A, \Pi \Sequent \Lambda$ برچسب می‌خورد؛ و اگر قاعده بر
!!{formula} نشاندار $\sFmla{\False}{!A}$ اعمال شود، برگ با
$\Pi \Sequent \Lambda, !A$ برچسب می‌خورد. هر جا تابلو قاعدهٔ
$\TRule{\True}{\star}$ را اعمال می‌کند، قاعدهٔ $\LeftR{\star}$ را
بر دنباله (با $!A$ همچون فرمول اصلی) اعمال کنید؛ و هر جا قاعدهٔ
$\TRule{\False}{\star}$ را اعمال می‌کند، $\RightR{\star}$ را به کار
برید. مقدمات قاعده همان دنباله‌هایی هستند که گره‌های متناظر افزوده‌شده
به تابلو را برچسب می‌زنند. وقتی تابلو با
$\sFmla{\True}{!A}$ و $\sFmla{\False}{!A}$ بسته می‌شود، برچسب گره برگ
پایانی دنباله‌ای به شکل
$!A, \Pi \Sequent \Lambda, !A$ خواهد بود. اگر شاخه با
$\sFmla{\True}{\lfalse}$ بسته شود، برچسب برگ پایانی دنباله‌ای با
$\lfalse$ در مقدم است؛ این نیز اصل \Log{G3c} است. برخی گره‌های پیاپی
تابلو دنبالهٔ یکسانی می‌گیرند؛ موارد تکراری را حذف کنید. برای نمونه،
تابلوی بالا به !!{proof} زیر ترجمه می‌شود:
\[\small
\Axiom$!A \fCenter !A, !B$
\RightLabel{\RightR{\lor}}
\UnaryInf$!A \fCenter !A \lor !B$
\Axiom$!A \fCenter !A, !C$
\RightLabel{\RightR{\lor}}
\UnaryInf$!A \fCenter !A \lor !C$
\RightLabel{\RightR{\land}}
\BinaryInf$!A \fCenter (!A \lor !B) \land (!A \lor !C)$
\Axiom$!B, !C \fCenter !A, !B$
\RightLabel{\RightR{\lor}}
\UnaryInf$!B, !C \fCenter !A \lor !B$
\RightLabel{\LeftR{\land}}
\UnaryInf$!B \land !C \fCenter !A \lor !B$
\Axiom$!B, !C \fCenter !A, !C$
\RightLabel{\RightR{\lor}}
\UnaryInf$!B, !C \fCenter !A \lor !C$
\RightLabel{\LeftR{\land}}
\UnaryInf$!B \land !C \fCenter !A \lor !C$
\RightLabel{\RightR{\land}}
\BinaryInf$!B \land !C \fCenter (!A \lor !B) \land
(!A \lor !C)$
\RightLabel{\LeftR{\lor}}
\BinaryInf$!A \lor (!B \land !C) \fCenter (!A \lor !B) \land
(!A \lor !C)$
\DisplayProof
\]

\end{document}
